\section*{Original introductory definitions}
First of all, we recall some basic definitions and notation. Given a compact Riemannian manifold $(N^{n+1},g)$ of dimension $n+1\geqslant 3$, with non-empty boundary $\partial N$, let $M^n$ be a codimension one properly embedded submanifold: we say that $M^n$ is a free boundary minimal hypersurface in $(N^{n+1},g)$ if it is a critical point for the $n$-dimensional area functional under variations satisfying the sole constraint that $\partial M\subset\partial N$ or, equivalently, if $M$ has zero mean curvature and meets the ambient boundary orthogonally. In order to avoid ambiguities, let us remark that properness is always understood here in the strong sense that $M \cap \partial N=\partial M$ (which is consistent with \cite{FL14, LZ16B, ACS17, ACS18b} among others).
    We let $\mathfrak{M}$ denote the class of smooth, connected, and properly embedded free boundary minimal hypersurfaces and consider the subclasses given by
     	\[
     		\mathfrak{M}(\Lambda,I) := \{ M \in \mathfrak{M} \ : \ \h^n(M) \leqslant \Lambda,\ \operatorname{index}(M) \leqslant I\},
     			\]
     			so with the additional requirements that area and Morse index be bounded from above (cf. \cite{Sha15}); and for $p$ an integer greater or equal than one
     \[
     			\mathfrak{M}_p(\Lambda,\mu) := \{ M \in \mathfrak{M} \ : \ \h^n(M) \leqslant \Lambda,\ \lambda_p \geqslant -\mu\},
     			\]
     			so with area bounded from above and $p^{th}$ eigenvalue $\lambda_p$ of the Jacobi operator bounded from below (cf. \cite{ACS15}).	
     			
     			
     		
     		In \cite{ACS17}, Ambrozio-Carlotto-Sharp proved a compactness theorem for $\mathfrak{M}_p(\Lambda,\mu)$ when the ambient dimension is less than eight: given a sequence of free boundary minimal hypersurfaces $\left\{M_k\right\}$ in $\mathfrak{M}_p(\Lambda,\mu)$, there is some smooth limit to which the sequence sub-converges smoothly and graphically (with finite multiplicity $m$) away from a discrete set ${\jepPubScript{Y}}$ on the limit, where curvature concentration occurs. These results apply, as a special case, to all classes of the form $\mathfrak{M}(\Lambda,I)$ (since this is clearly the same as $\mathfrak{M}_{I+1}(\Lambda,0)$). 
     			A significant part of the present article is aimed at an accurate description of the local picture around those points of bad convergence, in analogy with the results obtained, for the closed case, by Buzano and Sharp in \cite{BS17}. We then specify these results to the three dimensional scenario, and derive various sorts of new geometric results based on these tools.	\\
     			
     			
     			
     		As a first step in this program, we quantify the lack of smooth compactness (via a blow-up argument) by a finite number of non-trivial, complete, properly embedded minimal hypersurfaces of finite total curvature $\Sigma^n\emb \R^{n+1}$ as in \cite{BS17}, except this time it is possible that only `half' of $\Sigma$ is captured at a boundary point of bad convergence. In order to formalize this picture, let us introduce some terminology:
     		we employ the word \textsl{bubble} to denote a complete, connected and properly embedded minimal hypersurface of finite total curvature in $\R^{n+1}$; we employ the word \textsl{half-bubble} to denote a complete, connected, properly embedded minimal hypersurface that is contained in a half-space and has (non-empty) free boundary with respect to the boundary of this half-space,  and has finite total curvature. In both cases, the total curvature is understood to be the integral of the $n$-th power of the length of the second fundamental form, the corresponding functional being denoted by ${\jepPubScript{A}}(\cdot)$. Furthermore, we shall say that a (half-)bubble is \textsl{non-trivial} if it is not flat, namely if it is not a (half-)hyperplane.
     		
     		
\setcounter{section}{1}\setcounter{thm}{10}
